\section{Objectifs}

\begin{itemize}
\item Utiliser le vocabulaire image, antécédent, variable et courbe représentative.
\item Exploiter une formule y = f(x).
\item Lire graphiquement des images et des antécédents.
\item Modéliser une situation par une fonction.
\end{itemize}

\section{Prérequis}

\begin{itemize}
\item repère
\item lecture graphique
\item calcul littéral simple
\end{itemize}

\section{Activité de découverte}

La hauteur d'eau dans une cuve dépend du temps. À chaque instant, on associe une hauteur : cette dépendance se décrit par une fonction.

\section{Vocabulaire}

\begin{itemize}
\item \textbf{fonction}
\item \textbf{variable}
\item \textbf{image}
\item \textbf{antécédent}
\item \textbf{ensemble de définition}
\item \textbf{courbe représentative}
\end{itemize}

\section{Cours}

\subsection{Définitions}

\begin{itemize}
\item Une fonction associe à chaque valeur de la variable au plus une image.
\item L'image de x par f se note f(x).
\item La courbe d'équation y = f(x) est l'ensemble des points de coordonnées (x ; f(x)).
\end{itemize}

\subsection{Propriétés}

\begin{itemize}
\item Un point M(x ; y) appartient à la courbe de f si et seulement si y = f(x).
\item Un antécédent de k par f est une valeur x telle que f(x) = k.
\item L'ensemble de définition contient les valeurs de x pour lesquelles f(x) existe.
\end{itemize}

\subsection{Méthodes}

\begin{enumerate}
\item Pour calculer une image, remplacer x par la valeur donnée dans la formule.
\item Pour tester un point, remplacer x dans f(x) et comparer avec l'ordonnée du point.
\item Pour lire un antécédent sur un graphique, partir de l'ordonnée et rejoindre la courbe puis l'axe des abscisses.
\end{enumerate}

\section{Exemples corrigés}

\begin{itemize}
\item Si f(x) = 2x + 3, alors f(4) = 11.
\item Le point A(2 ; 7) appartient à la courbe de f(x)=2x+3 car f(2)=7.
\end{itemize}

\coursefigure{/images/mathematiques/latex-migration/2nde-fonctions-generalites-figure-1.svg}{Courbe affine avec lecture de l'image de 2.}{Fonctions : langage, courbes et modélisation : représentation schématique à commenter avant les exercices.}

\section{Algorithmique en Python}

\subsection{Une fonction Python comme programme de calcul}

\begin{verbatim}
def f(x):
    return 2*x + 3

for x in [-1, 0, 2, 4]:
    print(x, f(x))
\end{verbatim}

La fonction Python reçoit un argument x et renvoie son image par le programme de calcul.

\section{Erreurs fréquentes}

\begin{itemize}
\item Confondre f(2), qui est un nombre, et f, qui est la fonction.
\item Dire qu'une courbe est une fonction.
\item Lire l'image sur l'axe des abscisses au lieu de l'axe des ordonnées.
\end{itemize}

\section{Formules et idées à retenir}

\begin{itemize}
\item M(x ; y) appartient à Cf équivaut à y = f(x).
\item f(a) est l'image de a.
\end{itemize}

\section{Synthèse}

Ce chapitre sert à installer des automatismes, mais aussi à choisir une représentation adaptée : texte, formule, tableau, figure, graphique ou programme. Les exercices associés reprennent les cinq niveaux attendus : automatismes, application directe, méthode guidée, raisonnement et synthèse.
